\chapter{\textbf{Eigenvalue problems: vibrations and buckling}}
\label{chap:eigen}

\section*{Outline}
\addcontentsline{toc}{section}{\texorpdfstring{\hskip14mm}{} Outline}

The two previous chapters solved $KU=F$: given a load, find the response.
This chapter asks a different question -- for which loads, or at which
frequencies, does the structure respond \emph{without} being driven? Both
answers come from an eigenvalue problem, and in \texttt{Cast3M} both are
reached by the same route.

\begin{align}
  \text{free vibration:} \qquad & \left[K - (2\pi f)^{2} M\right] X = 0
  \label{eq:eigen:vibr} \\[1mm]
  \text{buckling:} \qquad & \left[K + \lambda\, K_{\sigma}\right] X = 0
  \label{eq:eigen:flam}
\end{align}

\cindex{eigenvalue problem}\cindex{vibration, free}\cindex{buckling}%
\cindex{natural frequency}\cindex{critical load}%
\noindent In the first, $M$ is the mass matrix and the eigenvalue is a
frequency. In the second, $K_{\sigma}$ is the \emph{geometric} stiffness
produced by a reference stress state, and the eigenvalue $\lambda$ is the
factor by which that reference load must be multiplied to make the structure
buckle. The mathematics is the same; only the second matrix changes.

We use the cylinder and the hyperbolic paraboloid of
section~\ref{sec:mesh:revolution}.

\section{Free vibration}

\subsection{Setting up the problem}

A vibration computation needs exactly one thing beyond a static one: a mass
matrix, which means a density in the material.

\begin{castemcom}
The model, the material and the stiffness are built exactly as in
chapter~\ref{chap:elasticity}; \op{mass[e]} \dindex{mass[e]}{mass matrix}
builds the mass matrix from the same model and material.

\sindex{coq4}\cindex{shell element}\cindex{shell thickness}%
\textbf{\texttt{'RHO'} is compulsory.} Without a density there is no mass
matrix, and the eigenvalue problem is meaningless. A shell also needs its
thickness, supplied through \op{cara} \dindex{cara[cteristiques]}{geometric characteristics} and concatenated to the material with \op{et} -- the single
most common reason a modal computation returns absurd frequencies.

The boundary conditions enter the stiffness exactly as before. Here the lower
rim of the cylinder is clamped. Note that \op{bloq 'DEPL' 'ROTA'} blocks both
translations and rotations, which is what an encastred shell edge needs.
\end{castemcom}
\begin{castemlines}
\begin{verbatim}
opti dime 3 elem qua4;

mod = cyl mode mecanique
          elastique isotrope coq4;

mat = mod mate 'YOUN' 2.1e11
               'NU'   0.3
               'RHO'  7800.;
car = mod cara 'EPAI' 0.002;
mat = mat et car;

KK = rigi mod mat;
MM = mass mod mat;

bl = bloq 'DEPL' 'ROTA' bas;

KT = KK et bl;
\end{verbatim}
\end{castemlines}

\subsection{Computing the modes}

\op{vibr[ation]} \dindex{vibr[ation]}{eigenvalue problem, vibration} solves
\eqref{eq:eigen:vibr}. It offers four algorithms, and which one you pick
matters more than it looks:

\sindex{vibr[ation], 'PROC'[HE]}\sindex{vibr[ation], 'SIMU'[LTANE]}%
\sindex{vibr[ation], 'INTE'[RVALLE]}\sindex{vibr[ation], 'IRAM'}%
\cindex{Lanczos algorithm}\cindex{Arnoldi algorithm}%
\begin{optable}{The four options of \texttt{vibr[ation]}}
'PROC'[HE] & the \texttt{n} modes nearest each frequency of a given list.
             Robust, and the one to use when you already know roughly where
             the modes are \\
'SIMU'[LTANE] & the \texttt{n} modes nearest one frequency, by Lanczos.
                Efficient when many modes are wanted -- \textbf{the sensible
                default} \\
'INTE'[RVALLE] & all modes in a frequency band, by bisection. The option to
                 use when the band matters rather than the count, but
                 generally far more expensive \\
'IRAM' & implicitly restarted Arnoldi, via \textsc{arpack}. The only option
         that accepts a damping matrix \\
\end{optable}

\begin{castemcom}
Here we ask for the ten lowest modes near zero, with \op{'SIMU'}. The
operator takes the stiffness and the mass matrix, in that order, and returns
a \op{table}.

\medskip
\textbf{The result structure is worth memorising}, because every
post-processing step goes through it. The table has an index \op{'MODES'},
itself a table indexed by mode number, and each entry holds:

\begin{itemize}[leftmargin=1.2em,itemsep=0pt,topsep=2pt]
\sindex{vibr[ation], 'MODES'}\sindex{'FREQUENCE'}%
\sindex{'DEFORMEE\_MODALE'}\sindex{'MASSE\_GENERALISEE'}%
\cindex{mode shape}\cindex{generalised mass}%
\item \texttt{'FREQUENCE'} -- the frequency, in Hz;
\item \texttt{'DEFORMEE\_MODALE'} -- the mode shape, a \op{chpoint};
\item \texttt{'MASSE\_GENERALISEE'} -- the generalised mass;
\item \texttt{'NUMERO\_MODE'} -- the mode number.
\end{itemize}

\noindent Frequencies come out in Hz, in whatever unit system the matrices
were built in -- so with SI data you get Hz, and with a density in
\texttt{g/mm}$^{3}$ you do not.
\end{castemcom}
\begin{castemlines}
\begin{verbatim}
nmod = 10;

res = vibr 'SIMU' 0. nmod KT MM;

* list the frequencies
lfreq = prog;
repeter bmod nmod;
  i  = &bmod;
  fi = res . 'MODES' . i . 'FREQUENCE';
  lfreq = lfreq et (prog fi);
  mess 'mode ' i ' : ' fi ' Hz';
fin bmod;
\end{verbatim}
\end{castemlines}

\begin{castemcom}
A mode shape is a \op{chpoint} like any displacement field, so it is plotted
with \op{defo} \dindex{defo[rmee]}{deformed shape} and \op{trac[er]}. Drawing
the undeformed shape with an amplification of \op{0.} and superposing it is
the usual way to make the mode legible.

\medskip
Remember that a mode shape has no physical amplitude -- it is normalised, here
by the generalised mass. Only its \emph{shape} and its frequency mean
anything; the numerical values of the displacements do not.
\end{castemcom}
\begin{castemlines}
\begin{verbatim}
repeter bdes nmod;
  i   = &bdes;
  di  = res . 'MODES' . i
            . 'DEFORMEE_MODALE';
  d0  = defo cyl di 0.  noir;
  dn  = defo cyl di 1.  vert;
  titr (chai 'mode ' i);
  trac cach (d0 et dn);
fin bdes;
\end{verbatim}
\end{castemlines}

\cindex{rigid body modes}\cindex{modal convergence}%
\noindent Two checks to run on any modal computation before trusting it.
First, count the zero frequencies. A free-free structure must return six of
them in 3D (three in 2D) -- its rigid body modes -- and a clamped one none.
A stray near-zero frequency means part of the mesh is not attached, very
often a \op{elim} that was forgotten, as on the cylinder seam.

\op{ense[mble]}, from section~\ref{sec:elas:ense}, answers this before you
run anything: it reports how many rigid body motions the restrained stiffness
still permits, so you know in advance how many zero frequencies to expect and
can tell a genuine one from a disconnected mesh. Second, refine the
mesh and watch the high modes: they converge far more slowly than the low
ones, and a mode whose wavelength spans only three or four elements is not
resolved.

\section{Buckling}

\subsection{The geometric stiffness}

Buckling asks at what load multiplier a structure loses stability. \cindex{geometric stiffness}\cindex{load multiplier}%
The
ingredient that makes it an eigenvalue problem is the \emph{geometric}
stiffness matrix $K_{\sigma}$: the change in stiffness produced by an existing
stress state. A membrane in tension is stiffer in bending; the same membrane
in compression is less stiff, and when the loss exactly cancels the elastic
stiffness, it buckles.

The computation therefore has two stages: a static computation under a
\emph{reference} load, and then the eigenvalue problem built on the stresses
it produced.

\begin{castemcom}
Before anything else, check the restraints with \op{ense[mble]}
(section~\ref{sec:elas:ense}). \op{flambage} needs a non-singular stiffness,
and an under-restrained model does not fail cleanly here -- it returns
critical multipliers that look plausible and mean nothing.

\medskip
Stage one is an ordinary elastic computation. The load is a reference load,
and its magnitude is arbitrary -- the eigenvalue $\lambda$ will scale it. Take
it equal to 1 in whatever unit is convenient and the critical load reads
directly off $\lambda$.

A linear eigenvalue analysis assumes a linear pre-buckling state, so the
stress is computed with the linear strain measure: \op{opti epsi lineaire}
and the \texttt{'LINE'} option of \op{sigm}, as in the shipped example
\texttt{flam1.dgibi}.

\op{ksigma} \dindex{ksigma}{geometric stiffness matrix} then builds
$K_{\sigma}$ from the model and the stress field.

\medskip
\textbf{Two traps in one line.} The keyword \texttt{'FLAM'} is
\emph{compulsory} if the matrix is to be used for buckling -- without it you
get a geometric stiffness built for a different purpose. And the geometric
characteristics \op{car} must be passed for the element families whose
stiffness depends on them -- the notice names thick shells, DST, beams and
pipes, and the shells used here carry a thickness, so pass it. It is an
optional \emph{positional} argument, so a missing \op{car} is not flagged as
a syntax error.
\end{castemcom}
\begin{castemlines}
\begin{verbatim}
* --- reference load: axial
*     compression of the cylinder
opti epsi lineaire;

fref = forc (0. 0. -1.) haut;

u0   = reso KT fref;
sig0 = sigm 'LINE' mod mat u0;

ksig = ksigma mod sig0 car 'FLAM';
\end{verbatim}
\end{castemlines}

\subsection{The \texttt{flambage} procedure}

\begin{castemcom}
\op{flambage} \sindex{flambage} is a \emph{procedure}, not an operator: it
takes a single \op{table} of inputs and returns a table of results, exactly
like \op{pasapas}.

\sindex{flambage, 'OBJM'}\sindex{flambage, 'CLIM'}\sindex{flambage, 'NMOD'}%
\sindex{flambage, 'LAM1'}\sindex{flambage, 'LAM2'}\sindex{flambage, 'SIG1'}%
\sindex{flambage, 'LAMB'}\sindex{flambage, 'DEPL'}%
The inputs are the model \texttt{'OBJM'}, the boundary conditions
\texttt{'CLIM'}, the number of modes \texttt{'NMOD'}, and the bracket
$[\texttt{'LAM1'},\texttt{'LAM2'}]$ within which the multiplier is sought.
The stress state is given either directly -- \texttt{'SIG1'} for the part that
scales with the multiplier, \texttt{'SIG0'} for a part that does not, plus
\texttt{'MATE'} and \texttt{'CARA'} -- or as the two matrices
\texttt{'RIGI'} and \texttt{'KSIG'} if you prefer to assemble them yourself.

\medskip
The bracket matters. Set \texttt{'LAM1'} slightly above zero, never to zero,
and \texttt{'LAM2'} well above the multiplier you expect; if the search
returns nothing, widen it before suspecting the model.

The result is a table indexed by mode number, each entry holding
\texttt{'LAMB'}, the critical multiplier, and \texttt{'DEPL'}, the buckling
mode shape.
\end{castemcom}
\begin{castemlines}
\begin{verbatim}
etab = table;
etab . 'OBJM' = mod;
etab . 'CLIM' = bl;
etab . 'MATE' = mat;
etab . 'CARA' = car;
etab . 'SIG1' = sig0;
etab . 'LAM1' = 0.001;
etab . 'LAM2' = 1.e6;
etab . 'NMOD' = 4;

stab = flambage etab;

repeter bflam 4;
  i  = &bflam;
  la = stab . i . 'LAMB';
  mess 'mode ' i
       ' critical factor ' la;

  mi = stab . i . 'DEPL';
  trac cach ((defo cyl mi 0. noir)
         et  (defo cyl mi 1. rouge));
fin bflam;
\end{verbatim}
\end{castemlines}

\noindent Because the reference load was a unit force, the critical load is
$\lambda$ itself. For the axially compressed cylinder, compare it with the
classical result
\[ N_{cr} = \frac{E\,t^{2}}{\sqrt{3(1-\nu^{2})}} \]
\cindex{imperfection sensitivity}\cindex{cylinder, axial buckling}%
per unit circumference. You will almost certainly find the computation
\emph{above} that value, and the discrepancy is the physics worth discussing:
the real cylinder is extremely sensitive to geometric imperfections, and the
linear eigenvalue analysis knows nothing about them. A buckling eigenvalue is
an upper bound, not a design load.

\subsection*{A warning about the name}

\dindex{flam}{combustion -- \emph{not} buckling}%
\cindex{flam, confusion with flambage}%
The procedure is \op{flambage}, spelled out. \op{flam} is a \emph{different
operator} -- it integrates hydrogen--oxygen combustion -- so the usual
four-character abbreviation does not apply here and silently calls the wrong
thing.

\section{Vibration under prestress}

\cindex{prestress, effect on frequencies}\cindex{vibration under prestress}%
The two problems combine. A stress state changes the stiffness, so it changes
the natural frequencies: a guitar string rises in pitch as it is tightened,
and a column in compression drops in frequency until, at the buckling load,
its first frequency reaches zero. That limit is the physical link between the
two halves of this chapter.

\begin{castemcom}
Adding the geometric stiffness to the elastic stiffness before calling
\op{vibr} is all that is needed. Here the cylinder carries an axial load
\op{pcharge} times the reference load.

Sweeping \op{pcharge} from zero up to the critical multiplier and plotting
the first frequency against it gives the curve that ends at zero -- the
clearest demonstration available that buckling is the point where a vibration
frequency vanishes.
\end{castemcom}
\begin{castemlines}
\begin{verbatim}
la1     = stab . 1 . 'LAMB';
pcharge = 0.5 * la1;

sigp = pcharge * sig0;
ksp  = ksigma mod sigp car 'FLAM';

KP = KK et ksp et bl;

resp = vibr 'SIMU' 0. 5 KP MM;

f1 = resp . 'MODES' . 1
          . 'FREQUENCE';
mess 'first frequency under load: '
     f1 ' Hz';
\end{verbatim}
\end{castemlines}

\section*{Summary}
\addcontentsline{toc}{section}{\texorpdfstring{\hskip14mm}{} Summary}

Both problems are built from the matrices of chapter~\ref{chap:elasticity}
plus one more:

\begin{center}
\begin{tabular}{@{}l@{\hspace{8mm}}l@{\hspace{8mm}}l@{}}
                   & \textbf{vibration} & \textbf{buckling} \\[1mm]
second matrix      & \op{mass} & \op{ksigma \ldots 'FLAM'} \\
needs              & \texttt{'RHO'} & a reference stress field \\
solved by          & \op{vibr} (operator) & \op{flambage} (procedure) \\
eigenvalue is      & a frequency, Hz & a load multiplier \\
result index       & \texttt{.'MODES'.i.'FREQUENCE'} & \texttt{.i.'LAMB'} \\
mode shape         & \texttt{.'DEFORMEE\_MODALE'} & \texttt{.'DEPL'} \\
\end{tabular}
\end{center}

\noindent Four things to carry away. Check the restraints with
\op{ense[mble]} first: both methods need to know how many rigid body motions
survive, and neither complains if the answer is wrong. A shell needs
\op{cara} with its thickness, passed everywhere the material is.
\op{ksigma} needs \texttt{'FLAM'}. And an eigenvalue computation is linear: it gives the
bifurcation point of the \emph{perfect} structure, which for
imperfection-sensitive shells can be well above the real collapse load.
Chapter~\ref{chap:largestrain} shows how to follow the post-buckling path
instead.

\section*{Exercises}
\addcontentsline{toc}{section}{\texorpdfstring{\hskip14mm}{} Exercises}

\begin{enumerate}
\item \textbf{Euler's column.} Mesh a cantilever beam with \op{coq2} elements
  and compute its first buckling load. Compare with
  $N_{cr}=\pi^{2}EI/(4L^{2})$. Then repeat with both ends pinned and with both
  clamped, and check that the critical loads are in the ratio $1:4:16$
  predicted by the effective lengths $2L$, $L$ and $L/2$. The shipped example
  \texttt{flam1.dgibi} is the model to start from.

\item \textbf{Modes of the cylinder.} Compute the twenty lowest modes of the
  clamped cylinder and sort them by the number of circumferential lobes.
  Which has the lowest frequency, and why is it not the one with the simplest
  shape? Refine the mesh and identify which modes have converged.

\item \textbf{The cooling tower.} Compute the first modes of the hyperbolic
  paraboloid of section~\ref{sec:mesh:revolution}, clamped at the base and
  free at the top. How does the lowest frequency change with the twist angle
  \op{alph}, at constant height and constant base radius? This is the
  structural reason for the shape.

\item \textbf{Frequency against load.} Sweep the axial load on the cylinder
  from zero to the critical value in ten steps, computing the first frequency
  at each. Plot $f_{1}^{2}$ against the load. The result should be close to a
  straight line hitting zero at the buckling load -- explain why, from
  \eqref{eq:eigen:vibr} and \eqref{eq:eigen:flam}.

\item \textbf{Imperfection sensitivity.} Add to the cylinder mesh a small
  geometric imperfection in the shape of its first buckling mode, with an
  amplitude of a fraction of the thickness, using \op{plus} as in
  section~\ref{sec:mesh:deform}. Recompute the buckling load. How fast does
  it fall as the amplitude grows?
\end{enumerate}
